% ============================================================================
% SINGLE-CLASS IMPLEMENTATION OF THE PRN FIX
%
% Idea: the paper's set-valued renegotiation functions are *already*
% conditional --- they take the counterpart's renegotiation function as an
% argument (Definition of \svrenegspacei: functions from
% \svrenegspacej x \allact to C(\allact)). So the structural constraint
% ("only add the PMP against counterparts who reciprocate") can live inside
% the renegotiation function itself. Consequences:
%   - Algorithm 2 (alg:csr) is UNCHANGED.
%   - The program spaces \fsspacei, \allprogreneg, and hence the statement
%     of the game \game(\allprog \cup \allprogreneg), are UNCHANGED
%     (a PMP-extended program is just a CSR program with a particular
%     renegotiation function).
%   - Only Definition 6 (def:pmp-extension) and Assumption 8
%     (assum:nopunish) change, plus the surrounding discussion and proof.
%
% This snippet contains: (0) a small piece of notation; (1) the revised
% Definition 6; (2) a well-definedness remark; (3) the revised Assumption 8;
% (4) replacement expository text for the passage introducing the
% assumptions (main.tex ll. 2879--2936); (5) the revised proof sketch;
% (6) notes on the required changes to the appendix proof, including a fix
% to a step that is stated incorrectly in the model's draft.
% ============================================================================


% ---------------------------------------------------------------------------
% (0) Notation: reference renegotiation outcome.
% (Place just before Definition 6. The appendix already defines a related
% object \bargdefault^{(k)} for ICSR; this is its one-shot analogue.)
% ---------------------------------------------------------------------------

For set-valued renegotiation functions~$\renegi, \renegj$ and an
outcome~$\fullact$, define the \termemph{reference renegotiation outcome}
\begin{align*}
\bargdefault(\renegi, \renegj, \fullact) =
     \begin{cases}
         \renegb\left(\renegi(\renegj, \fullact) \cap \renegj(\renegi, \fullact)\right),
         & \text{if } \renegi(\renegj, \fullact) \cap \renegj(\renegi, \fullact) \neq \emptyset; \\
         \fullact, & \text{otherwise.}
    \end{cases}
\end{align*}

% ---------------------------------------------------------------------------
% (1) Revised Definition 6.
% ---------------------------------------------------------------------------

\begin{definition}\label{def:pmp-extension}
Define the \termemph{PMP-extension} operator
$\mathcal{E}_i : \svrenegspacei \to \svrenegspacei$ as follows.
For any~$\renegi \in \svrenegspacei$, the PMP-extension
$\newrenegi = \mathcal{E}_i(\renegi)$ is given, for
all~$\rho^{j} \in \svrenegspacej$ and~$\fullact \in \allact$, by
\begin{align*}
    \newrenegi(\rho^{j}, \fullact) =
    \begin{cases}
        \renegi(\renegj, \fullact)
        \, \cup \,
        \yudproji\!\left(\bargdefault(\renegi, \renegj, \fullact)\right),
        & \text{if } \rho^{j} = \mathcal{E}_j(\renegj)
        \text{ for some } \renegj \in \svrenegspacej; \\[2pt]
        \renegi(\rho^{j}, \fullact),
        & \text{otherwise.}
    \end{cases}
\end{align*}
The PMP-extension of a CSR program~$\progi \in \fsspaceireni$ is the
program~$\newprogi \in \fsspaceinewreni$ with
$\defprog{\newprogi} = \defprogi$.
\end{definition}

% Consequence worth stating in-text right after the definition:
Notice that if both players use PMP-extensions, the agreement set computed
by Algorithm~\ref{alg:csr} is
\begin{align*}
    \left[\renegone(\renegtwo, \deffullact) \cup \yudprojone(\fullact^{\mathrm{def}})\right]
    \cap
    \left[\renegtwo(\renegone, \deffullact) \cup \yudprojtwo(\fullact^{\mathrm{def}})\right],
    \quad \text{where } \fullact^{\mathrm{def}} = \bargdefault(\renegone, \renegtwo, \deffullact),
\end{align*}
while against any counterpart whose renegotiation function is not a
PMP-extension, $\newrenegi$ returns exactly the same sets as~$\renegi$.
Thus a counterpart cannot gain from responding to the PMP-extension with a
more restrictive renegotiation set:\ against any such counterpart, the
PMP points are simply not offered, and the extension is behaviorally
indistinguishable from~$\renegi$.

% ---------------------------------------------------------------------------
% (2) Well-definedness remark (footnote or short remark after Def. 6).
% ---------------------------------------------------------------------------

% Suggested footnote text:
\footnote{For this definition to be well-posed, PMP-extensions must be
identified \textit{intensionally}:\ we take $\mathcal{E}_j(\renegj)$ to be
a canonical program implementing the two cases above, with the
base~$\renegj$ recoverable from its source code, so that
``$\rho^{j} = \mathcal{E}_j(\renegj)$ for some~$\renegj$'' is a syntactic
condition and the base is unique. (Identified extensionally, the
definition would be circular --- the first case quantifies over the image
of the operator being defined --- and the base need not be recoverable.)
This is the same device already used in Algorithms~\ref{alg:renegotiation}
and~\ref{alg:csr}, whose membership checks
``$\progj \in \fsspacej(\renegj)$ for some~$\renegj$'' are conditions on
the structure of the counterpart's code. Note also
that~$\mathcal{E}_i$ indeed maps~$\svrenegspacei$ to itself:\ since
every point of~$\renegi(\renegj, \fullact)$ weakly Pareto-improves
on~$\fullact$, so does~$\bargdefault(\renegi, \renegj, \fullact)$, and
hence so does every point
of~$\yudproji(\bargdefault(\renegi, \renegj, \fullact))$.}

% ---------------------------------------------------------------------------
% (3) Revised Assumption 8.
% Part (i) is unchanged. Part (ii) is replaced by the exact analogue of
% Assumption 4 / part (i): no program responds differently to a
% PMP-extension than to its base --- justified because the extension
% responds *identically* to any function that is not itself an extension.
% ---------------------------------------------------------------------------

\begin{assumption}
     We say that
     players
     with beliefs $\fullprior$
     \textbf{are (i) certain that CSR won't be punished and (ii) certain
     that PMP-extension won't be punished}
     if the following hold:
     \begin{enumerate}[(i)]
      \item Suppose either $\progi$ is in a subjective equilibrium
     of~$\game(\allprog \cup \allprogreneg)$,
     or $\progi$ is in the support of $\priorji$.
      Suppose $\progi \notin \fsspacei$.
         Then for any $\progj \in  \fsspacej$,
         we have
         $\progi(\progj) = \progi(\defprogj)$.
         \label{subassum:csrwlog}
         \item Suppose either $\progj \in \fsspacej(\rho^{j})$ is in a
     subjective equilibrium of~$\game(\allprog \cup \allprogreneg)$, or
     $\progj$ is in the support of~$\priorij$, where~$\rho^{j}$ is not the
     PMP-extension of any set-valued renegotiation function.
     Then for any~$\renegi \in \svrenegspacei$ with
     PMP-extension~$\newrenegi = \mathcal{E}_i(\renegi)$ and
     all~$\fullact$, we have
     $\rho^{j}(\newrenegi, \fullact) = \rho^{j}(\renegi, \fullact)$.
     \label{subassum:extension}
     \end{enumerate}
      \label{assum:nopunish}
\end{assumption}

% ---------------------------------------------------------------------------
% (4) Replacement expository text (for main.tex ll. 2879--2936, the passage
% introducing the two assumptions; adapt the lead-in as needed).
% ---------------------------------------------------------------------------

This result requires two assumptions on players' beliefs and the structure
of programs used in subjective equilibrium
(Assumptions~\ref{assum:nopunish}\ref{subassum:csrwlog}
and~\ref{assum:nopunish}\ref{subassum:extension}), both instances of the
same non-punishment principle as
Assumption~\ref{assum:rennopun} of Proposition~\ref{prop:simultaneous}:\
players do not respond differently to programs that do not respond
differently to them.
\begin{enumerate}
    \item Assumption~\ref{assum:nopunish}\ref{subassum:csrwlog}
is equivalent to
Assumption~\ref{assum:rennopun}
applied to CSR programs
rather than renegotiation programs:\ For any program
used in subjective equilibrium
or in the support of a player's
beliefs,
if that program never renegotiates,
it responds identically
to counterpart CSR programs
as to their defaults.
\item Assumption~\ref{assum:nopunish}\ref{subassum:extension} is the same
principle applied to renegotiation functions:\ no renegotiation function
used in subjective equilibrium or in the support of a player's beliefs
responds differently to the PMP-extension~$\newrenegi$ than to its
base~$\renegi$, unless it is itself a PMP-extension (in which case its
response is fixed by Definition~\ref{def:pmp-extension}).
This assumption is justified by the conditional structure of the
PMP-extension:\ by construction, $\newrenegi$ returns \textit{exactly the
same renegotiation sets} as~$\renegi$ against any function that is not a
PMP-extension. A counterpart therefore faces an identical decision problem
against~$\newrenegi$ as against~$\renegi$, and in particular cannot
benefit from restricting their renegotiation set in response to the
extension --- the PMP points they might hope to extract are only added
against counterparts who reciprocate.
(This is where the conditionality is essential. If player~$i$ added
their PMP \textit{unconditionally}, then~$\newrenegi$ would behave
differently toward~$j$ than~$\renegi$ does, and~$j$ might strictly prefer
a more restrictive renegotiation set against it --- for instance, to
induce the \selfunction{} function to choose among the added points. The
analogous assumption would then amount to assuming away a genuine
incentive, rather than ruling out gratuitous punishment.)
\end{enumerate}

% ---------------------------------------------------------------------------
% (5) Revised proof sketch for Theorem 3 (replaces main.tex ll. 3279--3329).
% ---------------------------------------------------------------------------

\begin{sketch}
Assumption~\ref{assum:nopunish}\ref{subassum:csrwlog}
implies that players always use CSR programs, so consider a subjective
equilibrium in CSR programs, and for each player~$i$ let~$\newprogi$ be
the PMP-extension of~$\progi$ (if~$\progi$'s renegotiation function is
already a PMP-extension, let~$\newprogi = \progi$).
Against any counterpart~$\progj$ in the support of~$i$'s beliefs whose
renegotiation function is not a PMP-extension, $\newrenegi$ returns the
same sets as~$\renegi$, and by
Assumption~\ref{assum:nopunish}\ref{subassum:extension}, $\progj$'s
renegotiation function responds identically to both; so the outcome is
unchanged.
Against a counterpart whose renegotiation function \textit{is} a
PMP-extension, the agreement set changes only by the addition of PMP
points of the reference renegotiation
outcome~$\fullact^{\mathrm{def}}$
--- all of which weakly Pareto-improve
on~$\fullact^{\mathrm{def}}$ --- so
because the \selfunction{} function is transitive, the new renegotiation
outcome is no worse for~$i$.
Therefore each player always weakly prefers to replace their program with
its PMP-extension, and the extended profile remains a subjective
equilibrium.
Finally, when all players use PMP-extensions, the agreement set contains
each player's PMP of~$\fullact^{\mathrm{def}}$, and no point that is worse
than the PMM for some player can be Pareto-efficient within it; so the
\selfunction{} function guarantees each player at least their PMM payoff.
\end{sketch}

% ---------------------------------------------------------------------------
% (6) NOTES ON THE APPENDIX PROOF (not paper text).
% ---------------------------------------------------------------------------
%
% (a) The reduction "wlog all programs are CSR" (main.tex ll. 3882--3899)
%     is unchanged: it only uses Assumption 8(i).
%
% (b) The case analysis for weak preference becomes:
%       - rho^j not a PMP-extension: identical outcomes, by 8(ii)
%         [replaces the old "S^PMP = empty / nonempty" split, which relied
%         on the old 8(ii)].
%       - rho^j = E_j(r^j): the agreement set is
%           [r^i(r^j,a) u PMP_i(a_def)] \cap [r^j(r^i,a) u PMP_j(a_def)],
%         a_def = ro(r^i, r^j, a). Relative to the reference agreement set
%         A = r^i(r^j,a) \cap r^j(r^i,a), the added points all lie in
%         PMP_i(a_def) u PMP_j(a_def), hence weakly Pareto-improve a_def.
%         Transitivity of sel then gives weak improvement for i.
%         NOTE: comparing i's OLD program against this counterpart uses
%         r^i evaluated at the argument E_j(r^j); Assumption 8(ii) (applied
%         to i's own equilibrium function r^i) equates this with
%         r^i(r^j, .). This is why 8(ii)'s scope covers equilibrium
%         programs as well as belief supports, mirroring 8(i).
%
% (c) The PMM-guarantee step needs one more line of care than in either
%     the original or the model's draft. It is NOT true in general that
%     every point of the extended agreement set that is below-PMM for some
%     player is Pareto-dominated by a point of
%     PMP_1(a_def) \cap PMP_2(a_def)  (such a point x could lie in
%     RN_1 \cap RN_2 with a very high payoff for the other player).
%     Correct route: let A' be the extended agreement set and x = sel(A').
%     (i) By transitivity, u(x) >= u(a_def).
%     (ii) If u_i(x) < PMM_i for some i, then x is in neither PMP_1 nor
%          PMP_2 (both require u_i >= max{PMM_i, u_i(a_def)} or equality),
%          so x lies in the reference agreement set A; since a_def = sel(A)
%          is Pareto-efficient in A and u(x) >= u(a_def), in fact
%          u(x) = u(a_def), so u_i(a_def) < PMM_i.
%     (iii) Any point m of PMP_1(a_def) \cap PMP_2(a_def) then satisfies
%          u_i(m) = max{PMM_i, u_i(a_def)} > u_i(x) and
%          u_j(m) = max{PMM_j, u_j(a_def)} >= u_j(x), and m is in A';
%          this contradicts Pareto-efficiency of x = sel(A') within A'.
%     CAVEAT (applies to the original proof and the model's draft alike):
%     step (iii) needs PMP_1(a_def) \cap PMP_2(a_def) to be NONEMPTY,
%     i.e., the payoff profile (max{PMM_1, u_1(a_def)},
%     max{PMM_2, u_2(a_def)}) to be feasible. Nonemptiness of each
%     PMP_i(a_def) separately (which is what the theorem's continuity
%     hypothesis gives) does not imply this. It holds, e.g., if the
%     feasible payoff set is comprehensive between the PMM and the
%     frontier; consider adding this as an explicit hypothesis or lemma.
%
% (d) Downstream edits elsewhere in the paper:
%     - Proposition 6 (prop:itspi): replace "Algorithm 2 with defaults p"
%       by "Algorithm 2 with renegotiation function E_i(r^i) and defaults
%       p"; the side condition PMP_i(a(p)) \subseteq RN_i(...) can be
%       dropped, since the extension now supplies exactly those points.
%     - The Remark after Theorem 3 (self-favoring extensions): the final
%       sentence ("it's plausible that Assumption 8(ii) would be
%       violated") should now say that a self-favoring extension does NOT
%       behave identically to its base against non-reciprocating
%       counterparts is false --- rather: the justification for 8(ii)
%       fails for self-favoring extensions, because the counterpart has a
%       genuine incentive to respond differently (the added points are
%       better for them and worse for the focal player), so no analogous
%       non-punishment assumption would be defensible.
%     - Proposition 5 (tightness): the last paragraph of its proof
%       verifies the old 8(ii) via "renegotiation sets independent of the
%       counterpart's function". Under the new 8(ii) the verification is
%       if anything easier (functions independent of the counterpart's
%       function trivially respond identically to E_i(r^i) and r^i), but
%       it should also be checked that the constructed programs remain
%       best responses when PMP-extensions are available as deviations;
%       since the constructed beliefs put probability 1 on non-extension
%       counterparts, extending changes nothing against them, so the
%       best-response argument is unaffected.
%     - The n-player statements in the appendix: E_i(r^i)(rho^{-i}, a)
%       adds PMP_i(ro(r^i, r^{-i}, a)) iff EVERY rho^j is a PMP-extension
%       (with r^{-i} the profile of bases), and behaves as r^i otherwise.
% ============================================================================
